\expandafter\ifx\csname CAPETINCLUDEMODE\endcsname\relax
\documentclass[11pt,a4paper]{article}
% preamble.tex — packages et commandes communes a toutes les fiches CAPET SII IE
\input{glyphtounicode}
\pdfgentounicode=1
\usepackage{cmap}
\usepackage[T1]{fontenc}
\usepackage[utf8]{inputenc}
\usepackage[french]{babel}
\usepackage[a4paper, top=1.6cm, bottom=1.6cm, left=1.4cm, right=1.4cm]{geometry}
\setlength{\headheight}{14.5pt}
\usepackage{amsmath, amssymb, amsfonts}
\usepackage{siunitx}
\sisetup{locale=FR, output-decimal-marker={,}, inter-unit-product=\ensuremath{{}\cdot{}}}
\usepackage{array, booktabs, multirow, tabularx}
\usepackage{xcolor}
\usepackage[most, breakable]{tcolorbox}
\usepackage{multicol}
\usepackage{enumitem}
\usepackage{mathtools}
\usepackage{lmodern}
\usepackage{microtype}
\usepackage{fancyhdr}
\usepackage{tikz}
\usetikzlibrary{positioning,arrows.meta}
\usepackage{lastpage}

% ── Couleurs par domaine ──────────────────────────────────────────────
\definecolor{colElec}{RGB}{26, 82, 148}       % Électrotechnique — bleu foncé
\definecolor{colElecLight}{RGB}{210, 225, 245} % fond bleu clair
\definecolor{colAuto}{RGB}{34, 120, 66}        % Automatique — vert
\definecolor{colMeca}{RGB}{185, 90, 20}        % Mécanique — orange
\definecolor{colTherm}{RGB}{160, 40, 40}       % Thermique — rouge
\definecolor{colQT}{RGB}{232, 244, 233}        % fond question type (vert très léger)
\definecolor{colJury}{RGB}{255, 243, 205}      % fond jury (jaune ambre)
\definecolor{colJuryFrame}{RGB}{200, 130, 0}   % bordure jury
\definecolor{colGrayLight}{RGB}{245, 245, 247} % fond gris neutre
\definecolor{colSeparator}{RGB}{180, 180, 180} % gris séparateur

% ── Style tcolorbox de section principale ────────────────────────────
\tcbset{
  ficheSection/.style={
    breakable,
    enhanced,
    colback=#1!7!white,
    colframe=#1,
    colbacktitle=#1,
    coltitle=white,
    fonttitle=\bfseries\normalsize,
    left=4pt, right=4pt, top=3pt, bottom=4pt,
    before skip=6pt, after skip=6pt,
  }
}

% ── Style encadré question type ──────────────────────────────────────
\tcbset{
  questiontype/.style={
    breakable,
    enhanced,
    colback=colQT,
    colframe=colElec!55!black,
    colbacktitle=colElec!55!black,
    coltitle=white,
    fonttitle=\bfseries\small,
    left=4pt, right=4pt, top=3pt, bottom=3pt,
    before skip=5pt, after skip=5pt,
  }
}

% ── Style encadré jury ───────────────────────────────────────────────
\tcbset{
  jury/.style={
    breakable,
    enhanced,
    colback=colJury,
    colframe=colJuryFrame,
    fonttitle=\bfseries\normalsize,
    left=4pt, right=4pt, top=3pt, bottom=3pt,
  }
}

% ── Mise en page ─────────────────────────────────────────────────────
\pagestyle{fancy}
\fancyhf{}
\renewcommand{\headrulewidth}{0.4pt}
\renewcommand{\footrulewidth}{0.3pt}

\setlist[itemize]{leftmargin=1.2em, itemsep=1pt, topsep=2pt, parsep=0pt}
\setlist[enumerate]{leftmargin=1.2em, itemsep=3pt, topsep=2pt, parsep=0pt}

\setlength{\parskip}{2pt}
\setlength{\columnsep}{10pt}

\def\CAPETBEGINFILE{\begin{document}}
\def\CAPETENDFILE{\end{document}}
\else
\def\CAPETBEGINFILE{}
\def\CAPETENDFILE{}
\fi

\fancyhead[L]{\textbf{\textcolor{colElec}{T11}} - Machine - courant continu (MCC)}
\fancyhead[R]{\small\textit{CAPET SII IE - Session 2026}}
\fancyfoot[C]{\small\thepage\,/\,\pageref{LastPage}}

\CAPETBEGINFILE

% -----------------------------------------------------------------------
%  BLOC EN-TeTE
% -----------------------------------------------------------------------

\begin{tcolorbox}[
  enhanced, breakable,
  colback=colElec,
  colframe=colElec,
  left=6pt, right=6pt, top=5pt, bottom=5pt,
  before skip=0pt, after skip=4pt
]
  {\color{white}\bfseries\large T11 \quad Machine - courant continu (MCC)}\\[2pt]
  {\color{colElec!30!white}\small
    \textbf{Domaine :} electrotechnique - Machines electriques \quad
    \textbf{Presence :} \textbf{8/13} sujets (2013-2026) - theme \emph{frequent}
  }\\[1pt]
  {\color{colElec!20!white}\footnotesize
    Sous-themes : equation electrique (FCEM, resistance d'induit) - couple electromagnetique - bilan de puissance - quadrants de fonctionnement - hacheur 1Q/4Q - demarrage rheostatique
  }
\end{tcolorbox}

% -----------------------------------------------------------------------
%  1. RAPPELS DE COURS
% -----------------------------------------------------------------------

\begin{tcolorbox}[
  ficheSection=colElec,
  title={1. Rappels de cours},
]

% -- 1.1 equations fondamentales ---------------------------------------
\textbf{1.1 equations fondamentales (excitation separee, flux $\phi$ constant)}

\smallskip
\begin{center}
\renewcommand{\arraystretch}{1.5}
\begin{tabular}{lll}
\toprule
\textbf{Grandeur} & \textbf{Expression} & \textbf{Convention} \\
\midrule
FCEM & $E = k\phi\Omega$ & $[k\phi] = \mathrm{V\!\cdot\!s\!\cdot\!rad^{-1}} = \mathrm{N\!\cdot\!m\!\cdot\!A^{-1}}$ \\[4pt]
equation electrique (moteur) & $U = E + R_a I_a$ & $U > E$, courant entrant \\[4pt]
equation electrique (generateur) & $U = E - R_a I_a$ & $E > U$, courant sortant \\[4pt]
Couple electromagnetique & $C_{em} = k\phi I_a$ & Nem \\[4pt]
Vitesse & $\Omega = \dfrac{U - R_a I_a}{k\phi}$ & rades$^{-1}$ \\
\bottomrule
\end{tabular}
\end{center}

\medskip
\hrule\vspace{4pt}

% -- 1.2 Bilan de puissance --------------------------------------------
\textbf{1.2 Bilan de puissance}

\smallskip
\[
  P_{abs} = U \cdot I_a = \underbrace{R_a I_a^2}_{P_{Joule}} + \underbrace{E \cdot I_a}_{P_{meca}}
  \qquad;\qquad
  P_{meca} = C_{em} \cdot \Omega
\]
Couple utile (en sortie d'arbre) : $C_u = C_{em} - C_{pertes}$ (pertes mecaniques, frottements).\\
Rendement : $\eta = P_{meca} / P_{abs} = E I_a / (U I_a)$.

\medskip
\hrule\vspace{4pt}

% -- 1.3 Quadrants de fonctionnement ----------------------------------
\textbf{1.3 Quadrants de fonctionnement}

\smallskip
\begin{center}
\renewcommand{\arraystretch}{1.25}
\begin{tabular}{llll}
\toprule
\textbf{Quadrant} & \textbf{Couple $C$} & \textbf{Vitesse $\Omega$} & \textbf{Mode} \\
\midrule
I  & $> 0$ & $> 0$ & Moteur direct \\
II & $< 0$ & $> 0$ & Generateur (frein recuperatif) \\
III & $< 0$ & $< 0$ & Moteur inverse \\
IV & $> 0$ & $< 0$ & Generateur inverse \\
\bottomrule
\end{tabular}
\end{center}

\medskip
\hrule\vspace{4pt}

% -- 1.4 Hacheur de commande -------------------------------------------
\textbf{1.4 Hacheur de commande de la MCC}

\smallskip
\begin{center}
\renewcommand{\arraystretch}{1.3}
\begin{tabular}{lll}
\toprule
\textbf{Hacheur} & \textbf{Tension moyenne} & \textbf{Quadrants} \\
\midrule
1 quadrant (serie) & $\langle U \rangle = \alpha \cdot E_{dc}$ & I seulement \\[4pt]
4 quadrants (pont en H bipolaire) & $\langle U \rangle = (2\alpha-1) \cdot E_{dc}$ & I, II, III, IV \\[4pt]
Demarrage - courant d'appel & $I_{demarr} = U/R_a$ ( $\Omega = 0$) & Limiter par rheostats ou hacheur \\
\bottomrule
\end{tabular}
\end{center}

\end{tcolorbox}

% -----------------------------------------------------------------------
%  2. FORMULAIRE SYNTHeTIQUE
% -----------------------------------------------------------------------

\begin{tcolorbox}[
  ficheSection=colElec,
  title={2. Formulaire synthetique},
]
\small
\begin{multicols}{2}
\renewcommand{\arraystretch}{1.45}
\begin{tabular}{@{}p{4.5cm}ll@{}}
\toprule
\textbf{Formule} & \textbf{Usage} & \textbf{U.} \\
\midrule
\multicolumn{3}{@{}l}{\footnotesize\textit{equations electriques}} \\
$U = E + R_a I_a$ & Moteur & V \\
$U = E - R_a I_a$ & Generateur & V \\
$E = k\phi\Omega$ & FCEM & V \\
$C_{em} = k\phi I_a$ & Couple em & Nem \\
$\Omega = (U - R_a I_a)/(k\phi)$ & Vitesse & rad/s \\
\midrule
\multicolumn{3}{@{}l}{\footnotesize\textit{Bilan de puissance}} \\
$P_{abs} = U I_a$ & electrique & W \\
$P_J = R_a I_a^2$ & Pertes Joule & W \\
$P_{meca} = E I_a = C_{em}\Omega$ & Mecanique & W \\
\bottomrule
\end{tabular}

\columnbreak

\renewcommand{\arraystretch}{1.45}
\begin{tabular}{@{}p{4.5cm}ll@{}}
\toprule
\textbf{Formule} & \textbf{Usage} & \textbf{U.} \\
\midrule
\multicolumn{3}{@{}l}{\footnotesize\textit{Commande par hacheur}} \\
$\langle U \rangle = \alpha E_{dc}$ & Hacheur 1Q & V \\
$\langle U \rangle = (2\alpha-1)E_{dc}$ & Pont en H & V \\
$\alpha = \frac{1}{2}(1 + \langle U \rangle/E_{dc})$ & Rapport cyc. 4Q & - \\
\midrule
\multicolumn{3}{@{}l}{\footnotesize\textit{Couple nominal}} \\
$\Omega_{nom} = 2\pi N_{nom}/60$ & Conversion & rad/s \\
$C_{nom} = P_{nom}/\Omega_{nom}$ & Couple nom. & Nem \\
$I_{demarr} = U/R_a$ & Courant 0 vitesse & A \\
\midrule
\multicolumn{3}{@{}l}{\footnotesize\textit{Donnees typiques}} \\
$E_{dc} = 207\,\mathrm{V}$ (PD2 + 230\,V) & Bus DC & V \\
$E_{dc} = 72\,\mathrm{V}$ (Urbanloop) & Bus DC & V \\
\bottomrule
\end{tabular}
\end{multicols}
\normalsize
\end{tcolorbox}

% -----------------------------------------------------------------------
%  3. QUESTIONS TYPES
% -----------------------------------------------------------------------

\begin{tcolorbox}[
  ficheSection=colElec,
  title={3. Questions types - extraites des sujets reels},
]

% -- QT1 -------------------------------------------------------------

% Mini-procedure de resolution type copie concours
\begin{tcolorbox}[questiontype,title={Procedure type copie concours (5 etapes)}]
\begin{enumerate}[leftmargin=1.5em]
  \item \textbf{Qualifier le modele} : hypotheses, domaine de validite, grandeurs d'entree/sortie.
  \item \textbf{Ecrire les lois} : equations de base avec symboles et unites SI.
  \item \textbf{Isoler l'inconnue} : expression litterale, puis application numerique tracee.
  \item \textbf{Controler l'ordre de grandeur} : coherence physique et unite finale.
  \item \textbf{Conclure concours} : phrase courte orientee decision technique.
\end{enumerate}
\end{tcolorbox}

% Visuel de methode : chemin de resolution CAPET
\begin{center}
\begin{tikzpicture}[font=\small, node distance=7mm and 6mm, >=Stealth]
  \node[draw, rounded corners, fill=colElec!8, inner sep=4pt] (n1) {Notions};
  \node[draw, rounded corners, fill=colElec!8, right=of n1, inner sep=4pt] (n2) {Methode};
  \node[draw, rounded corners, fill=colElec!8, right=of n2, inner sep=4pt] (n3) {Application annale};
  \node[draw, rounded corners, fill=colElec!8, right=of n3, inner sep=4pt] (n4) {Verification jury};
  \draw[->, thick, colElec] (n1) -- (n2);
  \draw[->, thick, colElec] (n2) -- (n3);
  \draw[->, thick, colElec] (n3) -- (n4);
\end{tikzpicture}
\end{center}
\begin{tcolorbox}[questiontype,
  title={QT1 - Demarrage rheostatique - courant constant (rame VMI) \hfill\normalfont\footnotesize E2 sujet 2019 (Q7-9)}
]
\textbf{enonce.}
4 moteurs CC en serie sont alimentes par $V_{alim} = 750\,\mathrm{V}$. Lors du demarrage, le courant est maintenu - $I = 380\,\mathrm{A}$ constant par des rheostats. La tension d'un moteur vaut $U_{moteur} = k_2 V_{train}^2$ avec $k_2 = 0{,}1\,\mathrm{V\!\cdot\!s^2\!\cdot\!m^{-2}}$. Donner la loi de commande $R_{th}(V_{train})$.

\tcbline

\textbf{Corrige.}

Avec 4 moteurs en serie et $R_a \approx 0$ :
\[
  V_{alim} = 4 U_{moteur} + R_{th} \cdot I
  \quad\Rightarrow\quad
  R_{th} = \frac{V_{alim} - 4 k_2 V_{train}^2}{I}
  = \frac{750 - 4 \times 0{,}1 \times V_{train}^2}{380}
\]
La resistance \textbf{diminue} quand le train accelere (la FCEM $E = k_2 V_{train}^2$ augmente et prend en charge la tension).\\
\textit{energie dissipee dans les rheostats = pertes par effet Joule lors du demarrage (toute l'energie stockee dans les inductances).}
\end{tcolorbox}

\vspace{4pt}

% -- QT2 -------------------------------------------------------------

\begin{tcolorbox}[questiontype,
  title={QT2 - Hacheur 4 quadrants - rapport cyclique nominal \hfill\normalfont\footnotesize E2 appliquee 2022 (banc d'essai automobile, Q35-37)}
]
\textbf{enonce.}
Un frein - courant de Foucault est alimente par un hacheur 4 quadrants (pont en H). $E_{dc} = 207\,\mathrm{V}$ (apres redresseur PD2). Tension nominale du frein : $\langle u_f \rangle = 192\,\mathrm{V}$.
\begin{enumerate}[label=(\alph*)]
  \item Justifier le choix d'un hacheur 4 quadrants.
  \item Calculer $\alpha$ en commande bipolaire.
\end{enumerate}

\tcbline

\textbf{Corrige.}

\textbf{(a)} Pour demagnetiser le frein, il faut inverser le courant (donc inverser la tension) - reversibilite en tension et courant necessaire - \textbf{hacheur 4 quadrants} (pont en H).

\textbf{(b)} Pont en H bipolaire : $\langle u_f \rangle = (2\alpha - 1) \cdot E_{dc}$
\[
  \alpha = \frac{1}{2}\left(1 + \frac{\langle u_f \rangle}{E_{dc}}\right)
  = \frac{1}{2}\left(1 + \frac{192}{207}\right)
  = \frac{1}{2}(1 + 0{,}927) \approx \mathbf{0{,}963}
\]
\end{tcolorbox}

\vspace{4pt}

% -- QT3 -------------------------------------------------------------

\begin{tcolorbox}[questiontype,
  title={QT3 - Couple nominal et verification de dimensionnement \hfill\normalfont\footnotesize E2 sujet 2019 (Q28)}
]
\textbf{enonce.}
Les moteurs Alsthom ont : $P_{nom} = 96\,\mathrm{kW}$, $N_{nom} = 1035\,\mathrm{tr/min}$. Calculer le couple nominal. Sachant que le couple thermique equivalent sur le trajet est $C_{th} \approx 557\,\mathrm{N\!\cdot\!m}$, conclure sur le dimensionnement.

\tcbline

\textbf{Corrige.}

\textbf{etape 1 - vitesse angulaire nominale :}
\[
  \Omega_{nom} = \frac{2\pi \times 1035}{60} = 108{,}4\,\mathrm{rad\!\cdot\!s^{-1}}
\]

\textbf{etape 2 - couple nominal :}
\[
  C_{nom} = \frac{P_{nom}}{\Omega_{nom}} = \frac{96\,000}{108{,}4} \approx \mathbf{885{,}6\,\mathrm{N\!\cdot\!m}}
\]

\textbf{Conclusion :} $C_{th} = 557\,\mathrm{N\!\cdot\!m} < C_{nom} = 885{,}6\,\mathrm{N\!\cdot\!m}$ - les moteurs sont \textbf{correctement dimensionnes}.\checkmark
\end{tcolorbox}

\end{tcolorbox}

% -----------------------------------------------------------------------
%  4. ATTENTES DU JURY
% -----------------------------------------------------------------------

\begin{tcolorbox}[
  jury,
  title={4. Attentes du jury - Erreurs recurrentes et points de vigilance},
]

\begin{enumerate}[label=\textbf{\arabic*.}, leftmargin=*, itemsep=5pt]

  \item \textbf{Convention moteur vs generateur.}
  En moteur : $U > E$, $U = E + R_a I_a$. En generateur : $E > U$, $E = U + R_a I_a$.
  La confusion est explicitement sanctionnee par le jury 2024.

  \item \textbf{Au demarrage, $E = 0$.}
  - $\Omega = 0$, la FCEM est nulle : $I_{demarr} = U/R_a$ est tres eleve.
  C'est pourquoi on insere des rheostats de demarrage ou on utilise un hacheur (demarrage progressif).

  \item \textbf{Couple electromagnetique $\neq$ couple utile.}
  $C_{em} = k\phi I_a$ est le couple developpe par les forces de Laplace.
  Le couple utile en sortie d'arbre est $C_u = C_{em} - C_{pertes\,mecaniques}$.

  \item \textbf{Unites du gain $k\phi$.}
  $[k\phi] = \mathrm{V\!\cdot\!s\!\cdot\!rad^{-1}} = \mathrm{N\!\cdot\!m\!\cdot\!A^{-1}}$ : le meme coefficient apparaet dans $E = k\phi\Omega$ et dans $C_{em} = k\phi I_a$.

  \item \textbf{Hacheur 4Q obligatoire pour la reversibilite.}
  Un hacheur 1Q (tension toujours positive, courant toujours positif) ne permet pas le freinage recuperatif ou l'inversion du sens de rotation. Justifier le choix du nombre de quadrants.

  \item \textbf{Bilan de puissance en 3 niveaux.}
  $P_{abs} = P_{Joule} + P_{meca}$. Presenter les trois termes avant tout calcul de rendement.

\end{enumerate}

\medskip
\begin{tcolorbox}[
  colback=colJuryFrame!12!colJury,
  colframe=colJuryFrame,
  left=4pt, right=4pt, top=2pt, bottom=2pt,
  fontupper=\small\itshape
]
Rapport jury 2022 : "Le jury apprecie les candidats qui maetrisent le bilan de puissance de la machine - courant continu et qui distinguent clairement la puissance electrique absorbee, les pertes Joule, et la puissance mecanique disponible."\\
Rapport jury 2023 : "Les candidats oublient souvent de justifier le choix du type de convertisseur (hacheur 1Q, 2Q ou 4Q) et se contentent de calculer le rapport cyclique sans argumenter sur la necessite de la reversibilite."
\end{tcolorbox}

\end{tcolorbox}

% -----------------------------------------------------------------------
%  PIED - PReSENCE DANS LES SUJETS
% -----------------------------------------------------------------------

\begin{tcolorbox}[
  enhanced, breakable,
  colback=colGrayLight,
  colframe=colSeparator,
  left=4pt, right=4pt, top=3pt, bottom=3pt,
  before skip=4pt, after skip=0pt,
  fonttitle=\small\bfseries,
  title={Presence dans les sujets analyses (2013-2026)},
]
\begin{center}
\footnotesize\renewcommand{\arraystretch}{1.15}
\begin{tabular}{lccccccccccccc}
\toprule
\textbf{Annee} & 13 & 14 & 15 & 16 & 17 & 18 & 19 & 20 & 21 & 22 & 23 & 24 & 25 \\
\midrule
E1 (commune)   & - & $\checkmark$ & - & - & $\checkmark$ & - & - & $\checkmark$ & - & $\checkmark$ & - & - & $\checkmark$ \\
E2 discipl. IE & $\checkmark$ & - & $\checkmark$ & - & - & $\checkmark$ & $\checkmark$ & - & - & $\checkmark$ & - & $\checkmark$ & - \\
\bottomrule
\end{tabular}
\end{center}
\smallskip
\footnotesize\textbf{Sous-themes frequents :}
equation electrique $U = E + R_a I_a$ (7/13) - couple electromagnetique (6/13) - bilan de puissance (5/13) - hacheur de commande 1Q/4Q (5/13) - demarrage rheostatique (3/13)

\medskip
\textbf{Ordres de grandeur (reperes rapides).}
\begin{itemize}[leftmargin=1.3em]
  \item Reseau BT FR : tension phase-neutre typique \SI{230}{\volt}, tension phase-phase typique \SI{400}{\volt}, frequence \SI{50}{\hertz}.
  \item Courants d'etude frequents en sujet : de quelques \si{\ampere} a quelques dizaines de \si{\ampere}.
  \item En conversion d'energie, un rendement attendu reste generalement entre 0{,}8 et 0{,}98 selon la technologie.
\end{itemize}

\medskip
\textbf{Synthese express (a refaire sans support).}
\begin{enumerate}[leftmargin=1.4em]
  \item Poser les hypotheses du modele et les conventions de signe.
  \item Ecrire 2 a 4 relations cles avec unites.
  \item Realiser un calcul numerique complet et verifier l'ordre de grandeur.
  \item Justifier le resultat par un argument technique de niveau CAPET.
\end{enumerate}

\end{tcolorbox}

\CAPETENDFILE









